\rulesection{22.0 Scenarios}

\rulesubsection{22.1 Set Up}

\textit{Shenandoah: Valley of Fire} has six scenarios. There is a Player Record Sheet for each scenario listing starting hexes or arrival turns for all units.

\subsection[22.2 The 1862 Scenarios]{22.2 The 1862 Scenarios (1, 2, and 3)}

The first three scenarios deal with Jackson's campaigns in the valley during 1862. The purpose of his campaign was to draw as many Union forces as possible away from McClellan and Richmond. To simulate this, the Union Panic Track is used for VPs and also to regulate Union forces. In these scenarios, Banks, Fremont, and Shields are all treated as supreme commanders. This is the only instance in the game where one side has more than one supreme commander.

\textbf{22.21 Victory Conditions.} During the Conclusion Phase of each turn, adjust the Union Panic Track using the listed shifts. Check the current box for VPs and record any earned that turn. At the end of the game, subtract the Union player's VP total from the Confederate player's total and compare the result:
\begin{itemize}
\item 10+ Confederate Overwhelming Victory
\item 8 or 9 Confederate Substantive Victory
\item 7 Confederate Marginal Victory
\item 5 or 6 Draw
\item 3 or 4 Union Marginal Victory
\item 2 Union Substantive Victory
\item 1 or less Union Overwhelming Victory
\end{itemize}

\textbf{22.22 Scenario 1. Kernstown to McDowell (21 March -- 8 May, 1862).}
\begin{itemize}
\item \textbf{22.221 Scenario Length:} 25 turns.
\item \textbf{22.222 Initiative:} Confederate has the initiative on 1--5; Union on 6.
\item \textbf{22.223 Union Panic Track:} Place marker in box 2.
\item \textbf{22.224 Special Rules:} The Confederate player may not destroy the supply base that begins at Mount Jackson. The Confederate player earns one VP if the base is moved to Staunton. Milroy, Schenk, and Detachment E may not move until Fremont enters on turn 4, though they may use reaction movement.
\end{itemize}

\textbf{22.23 Scenario 2. McDowell to Harper's Ferry (8 May -- 29 May, 1862).}
\begin{itemize}
\item \textbf{22.231 Scenario Length:} 11 turns.
\item \textbf{22.232 Initiative:} Confederate has the initiative on 1--5; Union on 6.
\item \textbf{22.233 Union Panic Track:} Place marker in box 3.
\item \textbf{22.234 Special Rules:} None.
\end{itemize}

\textbf{22.24 Scenario 3. Harper's Ferry to Port Republic (29 May -- 9 June, 1862).}
\begin{itemize}
\item \textbf{22.241 Scenario Length:} 6 turns.
\item \textbf{22.242 Initiative:} Confederate has the initiative on 1--5; Union on 6.
\item \textbf{22.243 Union Panic Track:} Place marker in box 4.
\item \textbf{22.244 Special Rules:} None.
\end{itemize}

\subsection[22.3 Scenario 4. Early in the Valley]{22.3 Scenario 4. Early in the Valley (26 June -- 6 July, 1864)}

\textbf{22.31 Victory Conditions.} The Confederates win if they have 8 or more VPs at the end of the game; otherwise the Union wins.

If Confederate combat units control the following at any time, they receive:
\begin{itemize}
\item Martinsburg: 1 VP
\item Harper's Ferry: 3 VP
\item Williamsport: 2 VP
\item Winchester: 1 VP
\end{itemize}

Each side receives 3 VPs for a major victory and 0 for a minor victory.

\textbf{22.32 Scenario Length:} 7 turns.

\textbf{22.33 Initiative:} Confederate has the initiative on 1--5; Union on 6.

\textbf{22.34 Special Rules:} On the first game turn only, the Confederate player may not roll on the Command Table, and the Union may use only reaction movement.

\subsection[22.4 Scenario 5. Sigel in the Valley]{22.4 Scenario 5. Sigel in the Valley (5 May -- 6 June, 1864)}

\textbf{22.41 Victory Conditions.} The Confederate player wins if he has devastated the hex containing Staunton (1243) by the end of the game. Otherwise the Union player wins.

\textbf{22.42 Scenario Length:} 16 turns.

\textbf{22.43 Initiative:} The Union player has the initiative on 1--4; the Confederate player on 5 or 6.

\textbf{22.44 Special Rules.} Beginning with Confederate movement on turn 8, Breckinridge and all Confederate infantry must move as quickly and directly as possible to Staunton and then be removed from play. Replacement points begin on turn 10. If Sigel loses a major battle, he is removed from play and replaced by Hunter; the SPs with Sigel are transferred to Hunter.

\subsection[22.5 Scenario 6. Sheridan in the Valley]{22.5 Scenario 6. Sheridan in the Valley (9 August -- 23 September, 1864)}

\textbf{22.51 Victory Conditions.} The Union player receives VPs for devastating these town hexes:
\begin{itemize}
\item Strasburg (3218): 1 VP
\item Woodstock (2624): 1 VP
\item New Market (2326): 1 VP
\item Luray (2827): 1 VP
\item Harrisonburg (1835): 2 VP
\item Cross Keys (1838): 2 VP
\item Honeyville (2532): 1 VP
\item Waynesboro (1644): 2 VP
\item Staunton (1243): 3 VP
\item Mechum River Station (2245): 3 VP
\item Charlottesville (2646): 3 VP
\end{itemize}

The Confederate player receives VPs for devastating:
\begin{itemize}
\item Romney (2110): 2 VP
\item Williamsport (4502): 3 VP
\item Harper's Ferry (4610): 5 VP
\end{itemize}

Both players receive 5 VPs for each major victory.

At the end of the game, subtract Confederate VPs from Union VPs and compare:
\begin{itemize}
\item 20+ Union Overwhelming Victory
\item 15 to 19 Union Substantive Victory
\item 10 to 14 Union Marginal Victory
\item 8 to 9 Draw
\item 5 to 7 Confederate Marginal Victory
\item 3 to 4 Confederate Substantive Victory
\item 2 or less Confederate Overwhelming Victory
\end{itemize}

\textbf{22.52 Scenario Length:} 23 turns.

\textbf{22.53 Initiative:} Union has the initiative on 1--4; Confederate on 5 or 6.

\textbf{22.54 Special Rules.} For the first 6 turns, Sheridan has a \texttt{-2} DRM on the Command Point Table. Replacement points begin on game turn 10.

\textbf{22.55 Scenario Extension.} To extend the scenario to 30 turns, have Rosser enter on turn 27 at Staunton and have Wilson replace Averill on turn 25. Add 6 points to each victory level before comparing earned victory points.
